\chapter{Huffman coding: short codes for common things}\label{ch:huffman}

\section{The intuition}

In Morse code, the most common letter in English, E, is a single dot, while rare letters such as Q take four symbols. The average message becomes much shorter than if every letter had the same length. \textbf{Huffman coding} builds the \emph{best possible} such code for a given set of symbol frequencies, with one restriction: each symbol gets a whole number of bits, and no code is the beginning of another code (a \emph{prefix-free} code), so the bit stream can be decoded without separators.

In our codec the ``symbols'' are the tokens from the last chapter: the 16 DC categories and the (run, size) bytes for AC values. A token such as ``EOB'' or ``run 0, size 1'' is extremely common, so it will get a 2--3 bit code; a rare token such as ``run 14, size 7'' may get 15 bits.

\section{Building the code: merge the two rarest}

\begin{enumerate}
\item Count how often each symbol occurs.
\item Take the two least frequent symbols, glue them under a new parent whose frequency is the sum.
\item Repeat with the new parent treated as a symbol, until one node (the root) remains.
\item The code of a symbol is the path from the root: left $=0$, right $=1$. Rare symbols hang deep in the tree and get long codes.
\end{enumerate}

\subsection{A textbook example (every number below is computed by the code)}

\begin{minipage}{0.47\textwidth}
\begin{tabular}{@{}cccc@{}}
\toprule
symbol & frequency & length & code\\ \midrule
\inputrows{snippets/huff_example}
\end{tabular}
\vspace{4pt}

Total: $\huffExBits$ bits for \huffExTotal\ symbols versus $\huffExFixedBits$ bits with a fixed 3-bit code. The theoretical lower bound (entropy) is $2.22$ bits per symbol and Huffman gives $2.24$.
\end{minipage}\hfill
\begin{minipage}{0.5\textwidth}
\centering
\begin{tikzpicture}[level distance=0.95cm,
  level 1/.style={sibling distance=2.6cm},level 2/.style={sibling distance=2.4cm},
  level 3/.style={sibling distance=1.9cm},level 4/.style={sibling distance=1.1cm},
  every node/.style={font=\scriptsize},
  inner/.style={circle,draw=accent,fill=plainbg,minimum size=0.55cm,inner sep=1pt},
  leaf/.style={rectangle,rounded corners=2pt,draw=soundfg,fill=soundbg,minimum size=0.5cm,inner sep=2pt}]
\node[inner] {100}
  child {node[leaf] {a\,45} edge from parent node[left,font=\tiny] {0}}
  child {node[inner] {55} edge from parent node[right,font=\tiny] {1}
    child {node[inner] {25} edge from parent node[left,font=\tiny] {0}
      child {node[leaf] {b\,13} edge from parent node[left,font=\tiny] {0}}
      child {node[leaf] {c\,12} edge from parent node[right,font=\tiny] {1}}}
    child {node[inner] {30} edge from parent node[right,font=\tiny] {1}
      child {node[leaf] {d\,16} edge from parent node[left,font=\tiny] {0}}
      child {node[inner] {14} edge from parent node[right,font=\tiny] {1}
        child {node[leaf] {e\,9} edge from parent node[left,font=\tiny] {0}}
        child {node[leaf] {f\,5} edge from parent node[right,font=\tiny] {1}}}}};
\end{tikzpicture}
\end{minipage}

\medskip
\textbf{Merge order:} f+e$=14$; c+b$=25$; 14+d$=30$; 25+30$=55$; a+55$=100$. The test \code{huffman\_textbook\_example} asserts exactly these code lengths.

\begin{plain}
Huffman is provably optimal among codes that give each symbol a whole number of bits: the average length $L$ always satisfies $H\le L<H+1$, where $H$ is the entropy (the information-theoretic minimum). The test \code{huffman\_kraft\_equality\_and\_optimal\_cost} checks this inequality on random distributions, and also that the code lengths satisfy $\sum 2^{-\ell_i}=1$ (the code ``uses up'' the whole binary tree, with no wasted codewords).
\end{plain}

\section{The implementation details that matter}

\subsection{Only the lengths are transmitted: canonical codes}

Sending the whole tree is wasteful. Instead, once the \emph{lengths} are known, the actual codes are assigned by a fixed rule so the decoder can recompute them: sort symbols by (length, symbol value), give the first one the all-zero code, and give each next symbol the previous code $+1$, shifting left whenever the length increases. In the example above: $b,c,d$ (length 3) get $100,101,110$; then the next code is $111$, shifted left once for length 4: $1110,1111$ for $e,f$.

\lstinputlisting{snippets/code_huff_canonical.cpp}

\subsection{What goes in the file}

For each of the two tables (DC and AC), the file stores 32 bytes ``how many codes have length 1, 2, \dots, 32'' followed by the used symbols in canonical order. That is the same idea as JPEG's own DHT segment. In our tests the DC table takes 34--42 bytes and the AC table 33--94 bytes.

\subsection{Decoding}

The decoder reads one bit at a time, builds up a code, and checks whether it falls into the range of codes of the current length. Because canonical codes of one length are consecutive numbers, the check is just ``is $\text{code}-\text{first}[\ell]<\text{count}[\ell]$''.
\lstinputlisting{snippets/code_huff_decode.cpp}
(This bit-by-bit loop is simple and correct but not fast; real decoders use lookup tables. See Chapter~\ref{ch:perf}.)

\subsection{Three engineering details}

\begin{description}
\item[Deterministic ties.] When two nodes have equal frequency, which is merged first? The code always merges the pair with the smallest (frequency, smallest symbol in the subtree), so the resulting lengths are unique, and the NumPy reference makes the same choice. This is what lets the two implementations produce byte-identical files.
\lstinputlisting{snippets/code_huff_merge.cpp}
\item[Length limit of 32.] Pathological frequency distributions (Fibonacci-like) can create a code longer than the 32 bits the bit writer handles. If the tree gets too deep, all frequencies are halved (rounded up) and the tree rebuilt, which flattens it. The test \code{huffman\_length\_limit\_is\_enforced} builds the worst case with 45 Fibonacci frequencies and checks the limit, the Kraft inequality, and a round trip of every symbol.
\item[A single used symbol.] A flat image uses only one DC and one AC symbol. A tree with one node would give a zero-length code; the code instead assigns length 1 so something is still written and read. The tests cover it.
\end{description}

\section{Why two tables, and why the amplitude bits are not Huffman coded}

DC symbols (categories of differences) and AC symbols (run/size bytes) have completely different statistics, so each gets its own table. The raw amplitude bits that follow a symbol are nearly uniformly distributed (any value in the category is about equally likely), so Huffman coding them would save almost nothing but cost a huge table.

\begin{sound}
``I compute symbol frequencies in a first pass, build two Huffman tables per image with a deterministic tie-break, convert them to canonical codes so only the code lengths need to be stored, and write the stream in a second pass. Canonical codes also make decoding a small range check per bit length. I verified optimality via the Kraft equality and the entropy bound in the unit tests.''
\end{sound}
